\documentclass[conference]{../../../../Conference-LaTeX-template_10-17-19/IEEEtran}
\usepackage{amsmath,amssymb,amsfonts,array,booktabs,tabularx,balance,url}
\newcolumntype{Y}{>{\raggedright\arraybackslash}X}

\title{Will Massive Spacecraft Ever Travel at the Speed of Light?\\
An Evidence-Bounded Research Program for Subluminal Relativistic Missions}

\author{\IEEEauthorblockN{Research Proposal}
\IEEEauthorblockA{Four-agent evidence synthesis -- physics, systems engineering, and human-operation design only}}

\begin{document}
\maketitle

\begin{abstract}
Travel at the speed of light is often framed as the decisive threshold for interstellar exploration. For a spacecraft with nonzero rest mass, however, special relativity makes exact attainment of the vacuum light speed $c$ a finite-energy impossibility: the Lorentz factor and ideal kinetic energy diverge as $v/c$ approaches one. That conclusion does not end the engineering problem. It relocates it to a more useful question: which subluminal velocity window is credible for a declared payload and mission after energy supply, acceleration, braking, propulsion optics or onboard power, interstellar-medium impacts, radiation, thermal rejection, communication, and crew requirements are considered jointly?

This report surveys a source-bounded literature on directed-energy propulsion, radiation-powered rockets, lightweight sails, interstellar gas and dust damage, space radiation, near-term interstellar probes, and propulsion null tests. It proposes the Relativistic Mission Envelope (RME), a provenance-preserving multi-objective model that labels a mission architecture feasible, infeasible, or incomplete. The central hypothesis is that a coupled model will produce a materially narrower and better calibrated subluminal feasible region than kinetic-energy or travel-time screening alone. The report supplies equations, variables, validation rules, conditional outcome branches, and a component-evidence pathway. It is a research proposal: no beam was fired, no spacecraft was flown, no material was exposed, and no human subject was studied. Its primary contribution is an auditable route from the inviolable light-speed boundary to decisive experiments and design choices for relativistic but subluminal robotic and, eventually, crewed missions.
\end{abstract}

\begin{IEEEkeywords}
interstellar travel, special relativity, directed energy, light sail, relativistic propulsion, interstellar medium, radiation shielding, mission design, human spaceflight
\end{IEEEkeywords}

\section{Introduction}
The prospect of travelling to another star motivates a deceptively simple question: will humanity ever travel at the speed of light? The phrase mixes two distinct targets. Light in vacuum propagates at $c$. A spacecraft is a massive system with payload, structure, energy source, heat rejection, sensors, and perhaps people. Under established special relativity, the latter cannot be accelerated to exactly $c$ by a finite amount of energy. No propulsion label--chemical, nuclear, antimatter, beam driven, or otherwise--changes that kinematic statement. The relevant scientific question is therefore not whether better engineering can overcome an infinite-energy requirement, but how close a mission may credibly approach $c$ before its coupled constraints dominate.

This reframing is not a retreat from ambitious exploration. Relativistic subluminal flight could radically change robotic mission times to nearby stars. Directed-energy studies derive relativistic sail motion while identifying diffraction and system parameters as terminal-speed constraints \cite{kulkarni2018}. Conceptual analyses of natural astrophysical sources likewise find that near-$c$ terminal speeds require idealized conditions and advanced materials, control, and interstellar-medium (ISM) handling \cite{lingam2020}. These are meaningful scientific starting points. They are not demonstrations that a macroscopic, crewed, braking-capable craft can reach $c$.

The central design error in public discussions is to quote a cruise velocity without reporting its mission definition. A gram-scale flyby sail, a robotic rendezvous vehicle, and a crewed vehicle have different mass, acceleration, shielding, life-support, reliability, and braking obligations. A calculation that omits braking cannot establish arrival and deployment capability. A calculation that treats ISM impacts as a postscript cannot establish vehicle survival. A calculation that converts proper-time reduction into a claim about human feasibility ignores radiation, acceleration, and closed-loop habitat requirements. These are linked constraints, not independent caveats.

This paper implements a four-agent research workflow. The Survey Agent freezes evidence and records the gaps that it supports. The Idea Agent selects the Relativistic Mission Envelope (RME), an integrated feasibility frontier rather than a single optimistic speed. The ExperimentDesign Agent defines prospective numerical studies and human-supervised component validation. The Author Agent translates those inputs into an IEEE-style plan without inventing a field result, propulsion breakthrough, or light-speed mission. The result is direct: exact light-speed travel for a massive craft is not an engineering objective supported by established physics; scientifically valuable subluminal relativistic missions remain a testable systems-engineering frontier.

\section{Physical Boundary and Evidence Survey}
\subsection{The massive-payload light-speed boundary}
Let $\beta=v/c$ and let $m$ be a craft's rest mass. The Lorentz factor and kinetic-energy floor are
\begin{align}
\gamma(\beta) &= \frac{1}{\sqrt{1-\beta^2}}, \label{eq:gamma}\\
K(\beta,m) &= \left[\gamma(\beta)-1\right]mc^2. \label{eq:kinetic-energy}
\end{align}
Equation \eqref{eq:kinetic-energy} is a lower bound on the energy acquired by the craft. It excludes conversion losses, propellant or beam-source losses, radiator mass, guidance, shielding, and braking. Since $\gamma$ diverges as $\beta$ tends to one, finite energy cannot produce $v=c$ for $m>0$. The result does not depend on current technology and should not be confused with a claim that no high-$\beta$ mission is possible. It specifies the limit that every such mission must respect.

The distinction matters for language. ``Near light speed'' may mean $0.2c$, $0.5c$, or an even larger fraction, each with radically different $\gamma$, kinetic energy, impact spectrum, and infrastructure. Time dilation can reduce elapsed proper time for a traveller, but it neither makes $\beta=1$ attainable nor removes the energy supplied in the launch frame. It also does not eliminate the requirement to decelerate if the mission is a rendezvous rather than a flyby.

\subsection{Propulsion evidence: possibility is conditional}
Directed-energy propulsion is the most useful near-term class for testing relativistic robotic concepts because energy can be produced away from a low-mass craft. Kulkarni \textit{et al.} derive a relativistic treatment of directed-energy acceleration and identify diffraction as a constraint on the terminal speed \cite{kulkarni2018}. Bae presents photon-propulsion concepts and emphasizes lengthy technological and economic development paths \cite{bae2012}. Heller \textit{et al.} show why extremely low areal density is central to solar-sail precursor missions, while their example remains a subrelativistic precursor rather than human transport \cite{heller2020}. These sources support a clear inference: sails may be a route to testing high-speed flight, but payload-to-area ratio, beam geometry, material response, and trajectory determine whether a particular velocity is credible.

Radiation-powered rockets expose a complementary trade-off. Fuzfa's relativistic analysis states that energy cost is the decisive challenge for radiation-powered interstellar missions, especially when payload mass is macroscopic \cite{fuzfa2019}. This is a systems result, not a dismissal of all propulsion research. It means that a credible study has to report where energy is generated, how momentum is transferred, how much is rejected as heat, and whether the stated architecture stops at its destination. Long's advanced fusion mission analysis illustrates a useful current comparison: modeled interstellar flyby and rendezvous cases operate in a subrelativistic range rather than at $c$ \cite{long2022}.

Two tempting shortcuts are excluded. First, beam power is not a free spacecraft resource: it requires an aperture, pointing, a source power budget, and a coupling model. Second, a speculative reactionless or anomalous-thrust assertion is not a design input until it survives calibrated reproduction. Neunzig \textit{et al.} report no net thrust in tested asymmetric infrared-laser-resonator configurations at their stated sensitivity \cite{neunzig2021}. Millis's assessment of breakthrough-propulsion research similarly reports null, unresolved, and follow-on outcomes rather than a verified way around known limits \cite{millis2005}. These results do not prove that every future discovery is impossible; they establish the evidentiary threshold for including a new mechanism in an interstellar mission model.

\subsection{Interstellar medium, radiation, and operations}
At relativistic velocity, the destination is not the only environment that matters. The forward ISM becomes a high-energy stream in the vehicle frame. Hoang \textit{et al.} model gas and dust interactions for a relativistic spacecraft and identify erosion, cratering, material modification, and shielding strategies as critical considerations \cite{hoang2017}. Drobny \textit{et al.} analyze implanted hydrogen and helium in materials and the resulting damage mechanisms for directed-energy relativistic craft \cite{droby2021}. These studies make a structural point: the mass assigned to shielding, thermal protection, and redundancy feeds back into the energy needed for acceleration and braking.

The ambient radiation environment is also dynamic. Guo \textit{et al.} review solar energetic particles, galactic cosmic rays, and related heliospheric-radiation limitations relevant to aerospace risk modeling \cite{guo2024}. Edelstein and Edelstein present a model-specific warning that neutral interstellar hydrogen becomes a severe passenger and instrument hazard at highly relativistic speed \cite{edl2012}. The appropriate interpretation is bounded. The work motivates a coupled dose, heat, and shield study; it does not define a universal safe-speed cutoff for all materials, trajectories, habitats, or future protection systems.

Near-term missions provide a scale check. The Interstellar Probe concept is designed for local interstellar and outer-heliosphere science with contemporary engineering constraints, not for relativistic cruise \cite{mcnutt2022}. Its value for this proposal is comparative: it demonstrates that ``interstellar'' is not synonymous with ``relativistic,'' and that telemetry, lifetime, trajectory, and system power remain central even when the target is far closer than another star.

\section{Research Gap and Selected Direction}
\subsection{Why a speed number is not a feasibility result}
The Survey identifies six linked gaps. The first is a mission-level speed optimum: kinetic energy and cruise time alone do not identify a useful $\beta$. The second is acceleration-braking closure. A flyby can avoid an arrival braking maneuver, but it cannot be presented as a settled orbital or landing mission. The third is the joint ISM, radiation, thermal, and shield-mass trade-off. The fourth is scale transition: a gram-scale sail result is not a direct argument for a human payload. The fifth is evidentiary discipline for claimed propulsion effects. The sixth is the human operational envelope, which joins proper acceleration, dose, habitat reliability, and governance.

The Idea Agent selects the \textit{Relativistic Mission Envelope}. RME is a framework that maps a declared architecture $\theta$ to a set of feasible, infeasible, or incomplete labels. Its novelty is not a claim of new physics. Its contribution is an evidence contract: no velocity point becomes a feasibility claim unless it declares mass, mission class, energy source, propulsion coupling, acceleration, braking, environmental scenario, protection mass, and operations class. The framework gives high-$\beta$ mission advocates a path to demonstrate progress while making missing constraints visible.

\subsection{Central hypothesis and falsifiers}
The central hypothesis is:
\begin{quote}
\textit{For a fixed mission distance and payload class, a coupled model of relativistic energy, propulsion, braking, ISM damage, radiation, and shielding yields a bounded subluminal feasible region that is materially narrower and better calibrated than energy-only or travel-time-only screening.}
\end{quote}

The hypothesis has three falsification routes. First, a validated architecture could lie outside the model's predicted infeasible region using the same declared payload and environmental conditions. Second, added modules could fail to improve out-of-sample classification over an energy-only baseline. Third, component data could demonstrate a protection or braking architecture whose mass and effectiveness invalidate the model's controlling high-$\beta$ trade-off. A lower-energy flyby does not falsify the claim about a coupled rendezvous envelope; those are different architectures.

\section{Experiment Design and Formal Framework}
\subsection{Mission-architecture state}
An experimental unit is a complete prospective mission vector
\begin{equation}
\theta=\{D,m_p,m_d,\beta,P,A,\sigma_s,a,B,\mathcal{E},\mathcal{O}\},
\label{eq:architecture-vector}
\end{equation}
where $D$ is destination distance, $m_p$ payload mass, $m_d$ dry mass, $P$ source power, $A$ optical or sail aperture, $\sigma_s$ shield areal density, $a$ proper acceleration, $B$ braking architecture, $\mathcal{E}$ an ISM/radiation environment scenario, and $\mathcal{O}$ the robotic or crewed operational class. The vector makes mass and mission class explicit before any calculation.

For an ideal interval of constant proper acceleration from rest, the proper-time and coordinate-distance relations are
\begin{align}
\tau_{\mathrm{acc}}(\beta,a) &= \frac{c}{a}\operatorname{atanh}(\beta), \label{eq:proper-time}\\
x_{\mathrm{acc}}(\beta,a) &= \frac{c^2}{a}\left[\gamma(\beta)-1\right]. \label{eq:accel-distance}
\end{align}
Equations \eqref{eq:proper-time} and \eqref{eq:accel-distance} are idealized kinematic components, not a propulsion design. A mission model adds a corresponding braking term when $B$ requires rendezvous. It must state the reference frame and cannot add travel times across incompatible frames.

The source energy is represented by an architecture-specific lower-bounded budget:
\begin{equation}
E_{\mathrm{source}} \geq \frac{K_{\mathrm{acc}}+K_{\mathrm{brake}}+E_{\mathrm{loss}}+E_{\mathrm{thermal}}}{\eta_{\mathrm{chain}}},
\label{eq:source-energy}
\end{equation}
where $\eta_{\mathrm{chain}}$ is the declared end-to-end efficiency and the numerator separates acceleration, braking, losses, and thermal-management energy. The equation deliberately prevents a writer from reporting Eq.~\eqref{eq:kinetic-energy} as the required electrical energy.

\begin{table*}[t]
\caption{Relativistic Mission Envelope modules and release rules. Every module is required for a finalized feasibility label.}
\label{tab:rme-modules}
\centering
\footnotesize
\begin{tabularx}{\textwidth}{@{}p{0.18\textwidth}Y Y Y@{}}
\toprule
\textbf{Module} & \textbf{Inputs and observables} & \textbf{Decision contribution} & \textbf{Validation / exclusion rule} \\
\midrule
Relativistic kinematics & $\beta$, mass, acceleration schedule, reference frame; $K$, $\tau$, distance & Enforces Eqs.~\eqref{eq:gamma}--\eqref{eq:accel-distance}; rules out a finite-energy $m>0$ light-speed point & Unit tests at low $\beta$, monotonic high-$\beta$ growth, independent equation implementation \\
Propulsion and optics & Source power, aperture, sail area, reflectivity, coupling efficiency, power duration & Estimates delivered momentum and energy; identifies beam, onboard, or hybrid assumptions & Compare to published directed-energy and sail regimes; do not import unreplicated anomalous thrust \\
Mission closure & Destination, flyby/rendezvous label, acceleration and braking paths, communication & Tests whether acceleration and braking fit in $D$ and whether arrival objective is stated & Any missing braking or mission-class field produces `INCOMPLETE` \\
ISM and shielding & Gas/dust scenario, frontal area, material layers, areal density, heat rejection & Produces damage, erosion, implantation, and thermal proxies; converts shield area to mass & Hold out environment cases; qualify a material only from human-supervised component tests \\
Radiation and operations & Ambient particle scenario, transformed forward exposure proxy, crew/robot class, dose and reliability limits & Separates robotic survival from crewed operational feasibility & No universal dose threshold; crew conclusions require domain-expert review \\
Uncertainty and decision & Parameter provenance, ranges, priors, confidence target, Pareto rank & Reports feasible, infeasible, or incomplete architecture set & Posterior-predictive and sensitivity checks; publish assumptions and exclusions \\
\bottomrule
\end{tabularx}
\end{table*}

\subsection{Damage, dose, and feasibility predicates}
The environment module models a forward-exposed area $A_f$, shield mass $m_s=A_f\sigma_s$, and a scenario-specific damage proxy $M_{\mathrm{dam}}(\theta)$. It also records a radiation/dose proxy $D_{\mathrm{rad}}(\theta)$. Neither proxy is called a material qualification or a medical result. Their purpose is to force mass and protection terms into the same architecture vector as velocity. The mass recursion is important: more shielding may reduce a damage proxy but increases the mass appearing in Eq.~\eqref{eq:kinetic-energy} and therefore the source budget.

RME assigns a feasibility state through
\begin{align}
C_E(\theta) &= \mathbb{I}\!\left[E_{\mathrm{source}}\leq E_{\max}\right], \nonumber\\
C_X(\theta) &= \mathbb{I}\!\left[x_{\mathrm{acc}}+x_{\mathrm{brake}}\leq D\right], \nonumber\\
C_M(\theta) &= \mathbb{I}\!\left[M_{\mathrm{dam}}\leq M_{\lim}\right], \nonumber\\
C_R(\theta) &= \mathbb{I}\!\left[D_{\mathrm{rad}}\leq D_{\lim}\right]\quad\text{for crewed missions}, \nonumber\\
F(\theta) &= C_E(\theta)C_X(\theta)C_M(\theta)C_R(\theta).
\label{eq:feasibility}
\end{align}
The binary portion of Eq.~\eqref{eq:feasibility} supports an auditable decision; the incomplete state prevents a nonreporting architecture from being quietly classified as viable. A probabilistic implementation reports a feasibility probability and the sensitivity of its label to each uncertain input, but it retains the same data-completeness gate.

\subsection{Prospective study phases}
The first phase is analytic verification. Independently implemented routines test dimensions, the low-$\beta$ approximation, monotonic energy growth, limiting behavior, and consistent reference-frame handling. The second phase sweeps three distinct mission classes: robotic flyby, robotic rendezvous, and crewed rendezvous. It computes Pareto fronts over travel time, source energy, payload fraction, distance consumed by acceleration/braking, damage proxy, and radiation proxy. It does not pool these classes into one generic ``interstellar vehicle.''

The third phase conducts ablations. Removing braking, shielding, ISM damage, radiation, or operations requirements measures exactly how a simplified analysis enlarges the apparent feasible region. Held-out environment and payload regimes test whether the coupled model improves decision prediction over an energy-only baseline. The fourth phase is a conditional component-evidence program: after facility, safety, and engineering approval, candidate shield coupons can be evaluated under relevant ion and thermal proxy conditions, and beam-sail optical/thermal predictions can be compared with calibrated bench measurements. These are future tests, not observations in this report.

\section{Mission Classes, Scaling, and Validation}
\subsection{Three mission classes, three different questions}
A useful RME study does not begin with an abstract ``starship.'' It begins with one of three mission classes. A robotic flyby minimizes arrival mass and can be scientifically attractive when an instrument package can acquire and relay high-rate measurements during a short encounter. It may use one-way acceleration and consequently avoids a destination braking requirement. Its principal question is whether an ultralight payload, sail, electronics package, and forward protection layer survive a specified transit and return a useful data volume. Directed-energy analysis and ISM-impact work make this the first class in which a high-beta demonstration is technically meaningful \cite{kulkarni2018,hoang2017}.

A robotic rendezvous has a more demanding contract. It must either carry a braking mechanism, exploit a destination environment, or rely on a previously installed braking infrastructure. The choice changes mass, acceleration distance, timing, navigation, and communication. An architecture that omits the term $x_{\mathrm{brake}}$ from Eq.~\eqref{eq:feasibility} may still be a flyby candidate; it is not evidence of an orbital or surface mission. This distinction also applies to multiple-stage trajectories. A launch stage and a braking stage can distribute the energy burden, but their stations, pointing, source power, and operational availability must enter the same mission vector.

A crewed rendezvous adds another layer rather than merely scaling mass. The acceleration schedule becomes an occupational exposure; radiation, dust, and equipment failure become life-support questions; and a habitat must retain air, water, food, thermal control, and repair capability over an interval that may still be long in the destination frame. Time dilation can change a crew's proper time, but it cannot make a missing braking architecture, radiation model, or closed-loop habitat disappear. RME deliberately permits a rigorous conclusion that an architecture is compelling for a robotic flyby while incomplete for human transport.

\begin{table}[t]
\caption{Mission-class separation required by the RME protocol.}
\label{tab:mission-classes}
\centering
\footnotesize
\begin{tabularx}{\columnwidth}{@{}p{0.21\columnwidth}Y Y@{}}
\toprule
\textbf{Class} & \textbf{Minimum closure} & \textbf{Failure mode that cannot be hidden} \\
\midrule
Robotic flyby & Payload, one-way acceleration, encounter geometry, telemetry, ISM survival & Reporting cruise time while omitting encounter duration, data return, or forward damage \\
Robotic rendezvous & All flyby fields plus braking energy/momentum and arrival operations & Calling a one-way sail trajectory an arrival or deployment mission \\
Crewed rendezvous & All rendezvous fields plus acceleration, radiation, habitat, medical, repair, and governance envelope & Scaling a robotic mass or survival proxy into a human-feasibility claim \\
\bottomrule
\end{tabularx}
\end{table}

\subsection{Scaling links that make high beta difficult}
The RME framework concentrates on the couplings that make a single attractive number fragile. First, the relativistic energy floor grows nonlinearly with $\beta$. Second, a protection layer adds mass proportional to its frontal area and areal density, so a design response to higher ISM exposure can increase the energy and braking burden. Third, reducing mass with a thinner sail or shield can expose materials to thermal, optical, and impact limits. Fourth, beamed systems exchange onboard propellant for remote infrastructure; diffraction and pointing become part of propulsion rather than secondary ground equipment. Fifth, power generated at the source, energy coupled to the sail, kinetic energy acquired by the vehicle, and heat rejected by the system are different accounting categories.

The correct output is therefore a frontier, not a winner declared from one variable. For an architecture family $\Theta$, RME seeks the nondominated set over travel time $T$, source energy $E_{\mathrm{source}}$, delivered payload $m_p$, protection mass $m_s$, and risk proxies. A point may have a shorter $T$ but require disproportionate energy, a much lower payload fraction, or a damage burden that cannot be supported by the assumed material. Such a point is useful for research prioritization even when it is infeasible: it identifies whether the next investment belongs in aperture scale, photon coupling, materials, shielding, braking, or a different mission objective.

The scale transition must be tested rather than assumed. A lightweight sail can exploit a high area-to-mass ratio; a crewed craft requires many subsystems whose area and mass scaling may follow different laws. The RME analysis therefore uses normalized mass categories, for example a microprobe, an instrumented robotic payload, and a protected crewed habitat, without pretending that a single numerical parameter represents all of them. Results are reported as families of assumptions and sensitivity intervals. No source in the frozen survey is allowed to transform a gram-scale trajectory into a human transport claim by simple multiplication.

\subsection{Verification hierarchy and parameter provenance}
The model is verified at three levels. At the equation level, unit tests reproduce the Newtonian limit of Eq.~\eqref{eq:kinetic-energy} at small $\beta$, check the divergence near one, and confirm that constant-proper-acceleration expressions use the declared frame. At the systems level, independent implementations compare energy, distance, and braking calculations over the same architecture vectors. At the evidence level, each input is tagged as a physical constant, a peer-reviewed estimate, a facility-calibrated measurement, an expert-reviewed assumption, or an unresolved placeholder. The final category forces the state `INCOMPLETE`.

Parameter provenance is as important as numerical precision. ISM density, dust size distribution, material loss, and radiation environment cannot be treated as universal constants for all directions and time windows. Each scenario receives a versioned environment card. The card includes the spatial path, distributional form, uncertainty interval, and source status. The same rule applies to beam aperture, source power, reflectivity, structural areal density, and braking effectiveness. This enables a reviewer to determine whether two attractive frontiers differ because of physics, an input range, or a missing system requirement.

Future component experiments are selected to break decision-relevant uncertainty, not to generate a generic technology demonstration. For protection concepts, decisive observables may include mass loss, subsurface gas accumulation, optical degradation, and heat transport under a clearly stated proxy environment. For a sail, they may include reflectivity, absorptivity, deformation, and pointing response under calibrated illumination. For an energy source, they may include duty cycle, phase stability, coupling efficiency, and aperture control. A component result advances a mission frontier only when it is linked to the parameter it constrains and includes a null or calibration path. This hierarchy directly addresses the Survey's requirement that unreplicated propulsion assertions not enter a baseline design \cite{neunzig2021,millis2005}.

\subsection{Uncertainty, ablation, and decision value}
RME uses a deliberately adversarial ablation plan. The energy-only model is the baseline. The braking module is added first, followed by ISM/thermal protection, ambient and transformed radiation, and then operational constraints. The change in the feasible set after each addition quantifies the cost of omitted physics. A model that is more complex but does not change mission decisions has failed to justify its complexity. A model that changes every decision but has unbounded parameter uncertainty has identified a measurement priority rather than a readiness claim.

The next measurement is selected by expected decision value. Suppose two architectures have similar travel time but differ in whether shielding mass or beam diffraction dominates their uncertainty. The priority is not the measurement that is easiest to make; it is the one most likely to change the Pareto ranking or feasibility state. This creates a disciplined research loop: survey evidence specifies gaps, the model exposes the gap's effect on a decision, and a future human-supervised test constrains that specific module. The loop can make progress toward robust subluminal missions without promising a solution to the light-speed boundary.
\section{Expected Outcome Branches}
The intended product is not a predetermined answer that a given speed is safe or impossible. If the coupled model removes high-$\beta$ points admitted by energy-only screening, and component evidence is compatible with its constraint modules, H-RME is supported. The resulting frontier identifies which quantity--source energy, braking distance, shield mass, thermal rejection, or operational envelope--dominates at each payload scale. That is a useful negative and positive result: it shows where engineering investment can alter the frontier and where it cannot.

If beam-sail, shield, or braking inputs expand a feasible region, the conclusion is conditional on those verified input ranges. If a robust architecture occupies a region labelled infeasible, the framework is revised and the responsible module is identified. If uncertainty makes two mission architectures indistinguishable, the next experiment is selected by expected information gain rather than by optimistic performance. If all crewed cases remain incomplete because no qualified dose, habitat, or long-duration acceleration input is available, the correct output is a robotic-only frontier, not an inferred human mission.

The model can also clarify how time dilation should be used. Traveller proper time is a reportable performance measure, but it never replaces the launch-frame energy, arrival-frame velocity, environmental exposure, or destination infrastructure. A high-$\beta$ flyby may be scientifically worthwhile even when a crewed rendezvous is not. Conversely, a subrelativistic architecture may be preferable when it carries a useful instrument suite, communicates reliably, and survives the environment. RME makes those comparisons explicit.

\section{Significance, Limits, and Human Review}
The scientific contribution is an integration discipline. It connects well-developed but often separated communities: relativistic dynamics, directed-energy optics, materials under high-velocity impacts, radiation environment modeling, spacecraft systems engineering, and human operations. It avoids the false choice between ``light speed is impossible'' and ``therefore interstellar flight is impossible.'' Exact $c$ remains excluded for massive craft under the governing framework; a succession of increasingly capable subluminal robotic missions can nevertheless test the technologies, models, and science cases needed for more ambitious travel.

The program has practical implications. A common mission vector permits like-for-like comparison between beam infrastructure, onboard fusion concepts, solar-sail precursors, and protection architectures. It forces a costed distinction between producing energy on the vehicle and producing it externally. It provides a transparent language for funding decisions: a concept can be promising at a component level, incomplete at a mission level, and still valuable as a targeted experiment. It also guards against using spectacular travel-time claims to obscure an omitted stopping maneuver or a nonexistent evidence path.

The principal limits are known. The literature inventory does not prove that all possible propulsion discoveries have been enumerated. Environmental distributions, long-duration material properties, biological dose-response thresholds, and future source-infrastructure performance can remain uncertain. The response is not to assume favorable values. RME reports their ranges, propagates uncertainty, and marks a claim incomplete when a necessary input is absent. Publisher and DOI landing-page verification of the frozen records remains a human-review requirement because the requested in-app browser was unavailable in the current sandbox.

Any physical implementation requires qualified human authority. High-power directed-energy systems, particle-beam or radiation facilities, vacuum systems, cryogenic systems, hazardous voltages, and flight hardware have facility-specific safety requirements. Any crewed inference requires independent radiation, medical, life-support, and ethics review. The present design authorizes none of these actions automatically.

\section{Conclusion}
Massive spacecraft will not travel at exactly the speed of light through finite engineering effort under established special relativity. This is not merely a current-technology limitation: the ideal kinetic-energy relation itself diverges at $\beta=1$. The important scientific frontier is therefore an evidence-bounded range of subluminal velocities, not a rhetorical promise of crossing $c$.

The proposed Relativistic Mission Envelope offers a way to study that frontier honestly and productively. It binds payload mass, energy source, acceleration, braking, ISM survival, radiation, protection mass, and operations to one declared architecture. It distinguishes flyby from rendezvous and robotic from crewed missions. It treats a new propulsion claim as a test obligation, not a blank check. The next meaningful result is a calibrated feasible frontier and a ranked list of measurements that would change it. Such a program can accelerate genuine interstellar capability while preserving the distinction between approaching the speed of light and claiming to attain it.

\appendices
\section{Provenance and Reporting Contract}
Every simulated or future measured architecture receives a persistent identifier, a source registry for each parameter, units, reference frame, payload and dry-mass breakdown, mission class, propulsion chain, environmental scenario, protection design, model version, and uncertainty representation. The report card lists the active constraints, the dominant source of variance, and all omitted fields. A result is not publishable as a feasibility statement if any required field is missing.

\section{Claim--Evidence Matrix}
\begin{table}[h]
\caption{Minimum evidence required for RME claims.}
\label{tab:claim-evidence}
\centering
\footnotesize
\begin{tabularx}{\columnwidth}{@{}p{0.27\columnwidth}Y@{}}
\toprule
\textbf{Claim} & \textbf{Minimum support and prohibited shortcut} \\
\midrule
Relativistic energy boundary & Equation, mass, frame, and source-energy chain; do not call kinetic energy the complete mission energy. \\
Relativistic propulsion performance & Declared beam/onboard architecture and independent model comparison; do not scale a gram sail directly to crewed mass. \\
Rendezvous capability & Acceleration and braking paths within distance; do not relabel a flyby as arrival. \\
Vehicle protection & Environment scenario plus damage, thermal, and shield-mass analysis; do not call a proxy a qualification test. \\
Crewed feasibility & Radiation, acceleration, habitat, and medical review inputs; do not infer it from robotic survival. \\
New propulsion mechanism & Calibrated null tests and independent replication; do not use an unreplicated claim in the baseline model. \\
\bottomrule
\end{tabularx}
\end{table}

\begin{thebibliography}{99}
\bibitem{kulkarni2018}
N. Kulkarni, P. Lubin, and Q. Zhang, ``Relativistic spacecraft propelled by directed energy,'' \textit{The Astronomical Journal}, vol. 155, no. 4, art. 155, 2018, doi: 10.3847/1538-3881/aaafd2.

\bibitem{hoang2017}
T. Hoang, A. Lazarian, B. Burkhart, and A. Loeb, ``The interaction of relativistic spacecrafts with the interstellar medium,'' \textit{The Astrophysical Journal}, vol. 837, no. 1, art. 5, 2017, doi: 10.3847/1538-4357/aa5da6.

\bibitem{fuzfa2019}
A. Fuzfa, ``Interstellar travels aboard radiation-powered rockets,'' \textit{Physical Review D}, vol. 99, art. 104081, 2019, doi: 10.1103/PhysRevD.99.104081.

\bibitem{lingam2020}
M. Lingam and A. Loeb, ``Propulsion of spacecraft to relativistic speeds using natural astrophysical sources,'' \textit{The Astrophysical Journal}, vol. 894, no. 1, art. 36, 2020, doi: 10.3847/1538-4357/ab7dc7.

\bibitem{bae2012}
Y. K. Bae, ``Prospective of photon propulsion for interstellar flight,'' \textit{Physics Procedia}, vol. 38, pp. 253--279, 2012, doi: 10.1016/j.phpro.2012.08.026.

\bibitem{edl2012}
W. A. Edelstein and A. Edelstein, ``Speed kills: Highly relativistic spaceflight would be fatal for passengers and instruments,'' \textit{Natural Science}, vol. 4, no. 10, pp. 749--754, 2012, doi: 10.4236/ns.2012.410099.

\bibitem{guo2024}
J. Guo \textit{et al.}, ``Particle radiation environment in the heliosphere: Status, limitations, and recommendations,'' \textit{Advances in Space Research}, 2024, doi: 10.1016/j.asr.2024.03.070.

\bibitem{droby2021}
J. Drobny \textit{et al.}, ``Damage to relativistic interstellar spacecraft by ISM impact gas accumulation,'' \textit{The Astrophysical Journal}, vol. 908, no. 2, art. 248, 2021, doi: 10.3847/1538-4357/abd4ec.

\bibitem{long2022}
K. F. Long, ``Development of the HeliosX mission analysis code for advanced ICF space propulsion,'' \textit{Acta Astronautica}, vol. 202, pp. 157--173, 2022, doi: 10.1016/j.actaastro.2022.10.022.

\bibitem{mcnutt2022}
R. L. McNutt \textit{et al.}, ``Interstellar probe - Destination: Universe!,'' \textit{Acta Astronautica}, vol. 196, pp. 13--28, 2022, doi: 10.1016/j.actaastro.2022.04.001.

\bibitem{neunzig2021}
O. Neunzig, M. Weikert, and M. Tajmar, ``Thrust measurements and evaluation of asymmetric infrared laser resonators for space propulsion,'' \textit{CEAS Space Journal}, vol. 14, no. 1, pp. 45--62, 2021, doi: 10.1007/s12567-021-00366-4.

\bibitem{heller2020}
R. Heller, G. Anglada-Escude, M. Hippke, and P. Kervella, ``Low-cost precursor of an interstellar mission,'' \textit{Astronomy \& Astrophysics}, vol. 641, art. A45, 2020, doi: 10.1051/0004-6361/202038687.

\bibitem{millis2005}
M. G. Millis, ``Assessing potential propulsion breakthroughs,'' \textit{Annals of the New York Academy of Sciences}, vol. 1065, no. 1, pp. 441--461, 2005, doi: 10.1196/annals.1370.023.
\end{thebibliography}

\end{document}
